\documentclass[11pt]{article}
\usepackage{fontspec}
\usepackage[french,english]{babel}
\usepackage{amsmath}
\usepackage{amssymb}
\usepackage{amsthm}
\usepackage{mathrsfs}
\usepackage{tikz-cd}
\usepackage{enumitem}

\newcommand{\printedpage}[1]{\marginpar{\footnotesize p.~#1}}
\newcommand{\msp}[1]{\marginpar{\footnotesize [#1]}}

% footnote marks printed as superscript (n) in the text, plain n before the note
\makeatletter
\renewcommand\@makefnmark{\hbox{\textsuperscript{(\@thefnmark)}}}
\renewcommand\@makefntext[1]{\parindent 1em\noindent\hbox to 1.8em{\hss\textsuperscript{\@thefnmark}\,}#1}
\makeatother

\setlength{\parindent}{0pt}
\setlength{\parskip}{0.6em}

\begin{document}

\printedpage{44}
% head: A. Grothendieck : Dérivateurs, Ch. XVI, Édition des Universités de Montpellier II et Paris VII
% foot: Version 18 mars 2019

\msp{page 82}
of the proof of 5.5 b).)

Let $x$ be an object (point) of $X$, $x_i$ its image in $X_i$. Let us consider (via my standard construction of categories of paths) the $\Omega^1(X,x)$, $\Omega^1(X_i,x_i)$. We find
\[ \Omega^1(X,x) \simeq \varinjlim_{I} \Omega^1(X_i,x_i), \]
and one checks that this is again an $\mathrm{L}_1$-$\pi_0$-adapted direct system. Iterating, we find for $n \in \mathbf{N}^*$
\[ \Omega^n(X,x) \simeq \varinjlim_{I} \Omega^n(X_i,x_i), \]
$\mathrm{L}_1$-$\pi_0$-adapted direct limit. Using 5.6, we conclude from this that for every $n$, the subset (let us suppose $I$ an ordered set, what the devil, that's enough!)
\[ I_n = \{ i \in I \mid \Omega^n(X_i,x_i) \ \text{ is 0-connected} \} \]
is cofinal in $I$, moreover it is stable under filtering $\varinjlim$ of cardinality $\leq \pi_0$. Thus % checked against PDF leaf 46: correct as printed; Theorem 5.4 explicitly uses pi_0 for the first infinite cardinal.
\[ I_\infty \overset{\text{def}}{=} \bigcap I_n \]
is also cofinal in $I$ (\textit{cf.} XVIII, 13.10.2, p. 302).

\msp{page 83}
Now $i \in I_\infty$ means that all the $\Omega^n(X_i,x_i)$ are 0-connected, which also means that $X_i$ is $\mathbf{W}_\infty$-aspheric (since it is already 0-connected) (\textit{cf.} XII, cor. p. 190). This completes the proof of theorem 5.4.

\section*{6\quad The condition Loc (5, 7 D$_{W_\omega}$) implies $W \subset \mathbf{W}_\infty$}

I am going to give here again, for the record, the proof of theorem 1 of XII:

\textbf{Theorem 6.1.} \textit{Let $W$ be a basic localizer satisfying} D$_{W_\omega}$ \textit{(where $W_\omega$ is the smallest of all basic localizers). Then we have $W = \mathrm{Fl}\,\mathrm{Cat}$,} or \textit{$W \subset \mathbf{W}_\infty$ (according as $W$ does not satisfy the connectedness axiom} Loc (5), \textit{or does satisfy it.})

I will give the proof assuming that $W$ satisfies Loc (5). For the proof of the fact that in the contrary case, we have $W = \mathrm{Fl}(\mathrm{Cat})$, I refer to XII, th. 2, p. 20. (See also the very close result, more general and weaker, recalled in prop. 2.7, page 12.)

\textbf{Lemma 6.2.} \textit{Let $W$, $W'$ be two basic localizers, $W$ satisfying} Loc (7 D$_{W_\omega}$). \textit{For $W \subset W'$ to hold, it is necessary and sufficient that} $\mathrm{Asph}_W \subset \mathrm{Asph}_{W'}$.

\msp{page 84}
It suffices of course to prove ``it suffices'', so let us suppose that every $W$-aspheric object is $W'$-aspheric, and let us prove that $W \subset W'$. So let $f : X \longrightarrow Y$ be in $W$, let us prove $f \in W'$. For this, I use the canonical factorization $f = pi$,
\[ X \xrightarrow{\ i \in W_\omega\ } \bar{X} \xrightarrow{\ p \in \mathrm{Parf}_{W_\omega}\ } Y. \]
\printedpage{45}
% head: A. Grothendieck : Dérivateurs, Ch. XVI, Édition des Universités de Montpellier II et Paris VII
% foot: Version 18 mars 2019
We therefore also have $i \in W$, and since $f = pi \in W$, we have $p \in W$, hence $p \in \mathrm{Parf}_{W_\omega} \cap W$. Now D$_{W_\omega}$ means that this set is contained in $W^{\mathrm{u}}$, so we have $p \in W^{\mathrm{u}}$, and in particular the fibers of $p$ are $W$-aspheric. They are therefore $W'$-aspheric, and since $p$ is $W_\omega$-perfect and a fortiori $W'$-perfect, it follows that $p \in W'$. We also have $i \in W_\omega \subset W'$, hence $i \in W'$, hence also $f = pi \in W'$, q.e.d.

Let us then return to the proof of 6.1. By the lemma, it suffices to prove that if $X$ in Cat is $W$-aspheric, it is $\mathbf{W}_\infty$-aspheric. But by the hypothesis Loc (5) we already know that $X$ is 0-connected, and choosing a point $x$ in $X$, its being $\mathbf{W}_\infty$-aspheric is expressed by the condition
\msp{page 85}
that the $\Omega^n(X,x)$ (for $n \geq 1$) are 0-connected (\textit{cf.} XII, cor. p. 190 already cited). For this, because of Loc (5) it suffices to prove that they are $W$-aspheric, and by induction it suffices to prove that $\Omega^1(X,x)$ is. Now by definition of $\Omega^1(X,x)$, this is the fiber at $x$ of an arrow in Cat
\[ E = \underline{\mathrm{Ch}}(X,x,-) \xrightarrow{\ p \in \mathrm{Parf}_{W_\omega}\ } X, \]
where the source $E$ is $W_\omega$-aspheric, a fortiori $W$-aspheric. Since $X$ is as well, we therefore have $p \in W$, hence $p \in \mathrm{Parf}_{W_\omega} \cap W \subset W^{\mathrm{u}}$, hence $p$ has $W$-aspheric fibers, and in particular $\Omega^1(X,x)$ is $W$-aspheric, q.e.d.

\textbf{Corollary 6.3.} (Characterization of $\mathbf{W}_\infty$.) \textit{$\mathbf{W}_\infty$ is the largest basic localizer satisfying} Loc (5, 7 D$_{W_\omega}$) \textit{- hence also the largest basic localizer satisfying} Loc (5, 7).

I think that we will find below an ``internal'' characterization of $\mathbf{W}_\infty$, via a ``descent axiom'' (very strong) Loc (10),
\msp{page 86}
as the unique basic localizer satisfying Loc (5, 7 D$_{W_\omega}$, 10) (or again Loc (5, 7, 10)).

6.4. \textit{Comment on the axioms for $W$-fibrations of the types} Loc (7 \ldots), \textit{and in particular on} Loc (7 D$_{W_\omega}$). \textit{Conjectural relation with the} ``simple homotopy types'' \textit{of} \textsc{Whitehead}. % [sic] du types -> du type (or des types), grammatical agreement

a) It seems to me that \textit{all} the basic localizers with which we have tacitly worked up to now (the $\mathbf{W}_n$ ($n \in \mathbf{N}^*$), homological and cohomological equivalence, rational equivalence \ldots) correspond to localizers $W \supset \mathbf{W}_\infty$. By theorem 6.1, it therefore follows that \textit{none} of them, apart from $\mathbf{W}_\infty$ itself, satisfies the very weakest version D$_{W_\omega}$ of axiom Loc (7 \ldots), with the sole exception of Loc (7 F). Thus, the localizers $W$ which satisfy it are with the exception of D$_{W_\omega}$ itself, of a quite new type - localizers which are strictly contained in $\mathbf{W}_\infty$.

b) This raises, of course, the question of the existence of basic localizers contained in $\mathbf{W}_\infty$ and
\msp{page 87}
distinct from $\mathbf{W}_\infty$. I have not even proved that $W_\omega \neq \mathbf{W}_\infty$, although I have not the slightest doubt that this is indeed the case. \texttt{[Cisinski has proved that} $W_\omega = \mathbf{W}_\infty$\texttt{.]} A fortiori, I do not know whether the smallest basic localizer satisfying Loc (7 D$_{W_\omega}$) is different from $\mathbf{W}_\infty$. I would be quite particularly interested in the smallest localizer $\tilde{\mathbf{W}}_\infty \subset \mathbf{W}_\infty$ satisfying all the stability axioms Loc (1 ter, 2, 3 bis, 4, 6, 7, 8) (\textbf{N.B.} it satisfies
\printedpage{46}
% head: A. Grothendieck : Dérivateurs, Ch. XVI, Édition des Universités de Montpellier II et Paris VII
% foot: Version 18 mars 2019
Loc (5) necessarily, since $\tilde{\mathbf{W}}_\infty \subset \mathbf{W}_\infty \subset \mathbf{W}_0$, moreover Loc (7) is equivalent, in the presence of Loc (8), to the very weak form Loc (7 D$_{W_\omega}$): $\mathrm{Parf}_{W_\omega} \cap W \subset W^{\mathrm{u}}$. Finally, one shows that Loc (4) is a consequence of the other axioms, \textit{cf.} below.) I now conjecture that $\tilde{\mathbf{W}}_\infty$ (which gives rise to a homotopy formalism just as complete as the everlasting ``weak equivalence'') is indeed different from $\mathbf{W}_\infty$ \texttt{[no by Cisinski]}. And I suspect
\msp{page 88}
that the category $\mathrm{Hot}_{\tilde{\mathbf{W}}_\infty}$ deserves to be called the category of \textit{``simple homotopy types''} (in the style of \textsc{Whitehead})\footnote[41]{Botched (?), \textit{cf.} comment on page 92.}. A first decisive test is to check that the principal cohomological invariant of \textsc{Whitehead} of a simple homotopy type, namely an element in a $K(\mathbf{Z}[\pi_1])$ (if I remember correctly\footnote[42]{no, it is a finer invariant \ldots}) - or something close to it, can already be defined in terms of the objects in $\mathrm{Hot}_{\bullet\,\tilde{\mathbf{W}}_\infty}$ (the pointed Hot). And secondly, if one associates to a finite ``simplicial mock-up'' (describing combinatorially a finite triangulation) the ordered set of its simplices, viewed as an object of Cat, and the corresponding object in $\mathrm{Hot}_{\tilde{\mathbf{W}}_\infty}$, the latter does not change\footnote[43]{This is obvious even for $W_\omega$, since the homotopy type \textit{in the strict sense does not change!} - No, there is confusion here between topological homotopy and combinatorial homotopy, \textit{cf.} p. 92 \ldots}, up to canonical isomorphism, when one performs on this ``polyhedron'' under consideration an ele\msp{page 89}mentary operation in the style of \textsc{Whitehead}. This moreover suggests trying to define such elementary operations in the whole of Cat, or at least in the subcategory of (small) ordered sets, in the hope of arriving (perhaps) at a more direct construction of the localizer relevant to simple homotopy type.

c) I end with a comment on the meaning of axiom Loc (7 D$_{W_\omega}$), as opposed to the weaker axiom Loc (7 F) which only ensures the existence of a ``structure with fibrations'' $(W, \mathrm{Fib}_W)$ on Cat. I also point out in passing that Loc (7 F), just like Loc (5) and like Loc (9), and unlike all the other axioms reviewed up to now, is \textit{not} a stability axiom: it asserts, not that certain arrows (under suitable hypotheses) are in $W$, but the existence of a factorization. Thus it is
\msp{page 90}
not at all clear that the intersection of two basic localizers satisfying Loc (7 F) satisfies it.

Loc (7 D$_{W_\omega}$) is exactly what is needed to ensure that the canonical factorization of an arrow $f : X \longrightarrow Y$ of Cat as $f = pi$, with $i \in W_\omega$, $p \in \mathrm{Parf}_{W_\omega}$, is also a factorization of the type stipulated by the factorization axiom Loc (7 F): we have (of course) $i \in W$, but also $p \in \mathrm{Fib}_W$ (since D$_{W_\omega}$ is equivalent to its strong form $\mathrm{Parf}_{W_\omega} \subset \mathrm{Fib}_W$). This also means that the canonical factorization of $f$ can be used to ``compute'' its homotopy fibers, not only at the points of $Y$, but also with respect to a base change morphism $Y' \longrightarrow Y$.

In the absence of this axiom,
\msp{page 91}
\printedpage{47}
% head: A. Grothendieck : Dérivateurs, Ch. XVI, Édition des Universités de Montpellier II et Paris VII
% foot: Version 18 mars 2019
assuming only Loc (7 F), we know that there exists a theory of $W$-homotopy fibers, but we have a priori no practical means of ``computing'' the said fibers, for lack of a manageable criterion sufficient to ensure $p \in \mathrm{Fib}_W$. I therefore have the impression that axiom Loc (7 F) is going to be of meager use. (It is true that the example of $\mathbf{W}_0$ ought to make me cautious! I would have to look at the case of the $\mathbf{W}_n$, whether they satisfy Loc (7 F) \ldots). To my mind, axiom Loc (7 D$_{W_\omega}$) is the exact ``bare minimum'' for being at ease in handling $W$-homotopy fibers. It is remarkable that this implies that one must have $W \subset \mathbf{W}_\infty$, and that this at the same time excludes \textit{all} the variants considered to date of the notion of homotopy type!

\msp{page 92}
6.5. Finally, the relation glimpsed with the ``simple homotopy types'' of \textsc{Whitehead} vanishes\footnote[44]{But perhaps not after all - the question remains open.}: indeed, the latter are strictly finer than ordinary homotopy types (\textit{strict}), and Whitehead's invariant is \textit{not} an invariant of ordinary homotopy type. It is true that when one considers the simplicial skeleton of a triangulated polyhedron as an object of Cat, ordinary (continuous) homotopy between two such polyhedra is interpreted, in Cat, not (it seems to me) by the existence of a homotopism between the corresponding objects of Cat, but by their being isomorphic in $\mathrm{Hot}_{\mathbf{W}_\infty} \overset{\text{def}}{=} \mathrm{Cat}(\mathbf{W}_\infty)^{-1}$. But the method of simplicial approximations of \textsc{Brouwer} shows that if the polyhedra $X$ and $Y$ are homotopic, then there exist simplicial subdivisions $X'$ and $Y'$ of $X$ and of $Y$ which are ``combinatorially homotopic'' (in a suitable sense). So I presume that the simplicial mock-ups of $X'$ and $Y'$ are then homotopic in Cat. It would follow that % [sic] plus fin -> plus fins, agreement with "ces derniers"
\msp{page 93}
$X$ and $Y$, I mean their mock-ups, are indeed homotopic in Cat, if it is true that a simplicial subdivision (barycentric subdivision would suffice, iterating \ldots) $X'$ of $X$ has a mock-up in Cat which is homotopic to it. The question raised here seems interesting to me, clearly one must find the answer to it. It already arises for any ordered set $I$, which we will suppose finite to simplify, taking the set (ordered by inclusion) $I' = \mathrm{sd}(I)$ of totally ordered subsets of $I$. We know that the topological realizations of $I$ and of $\mathrm{sd}(I)$ are canonically isomorphic, but are $I$ and $\mathrm{sd}(I)$ homotopic in Cat? We do have a canonical increasing map % [sic] ensemble ordonnée -> ensemble ordonné, agreement
\[ f : \mathrm{sd}(I) \longrightarrow I, \]
which associates to every totally ordered subset of $I$ its greatest element. But is there an increasing map in the opposite direction, which is a homotopism inverse to the preceding one up to homotopy? I have just checked in the case

\msp{page 94} missing

\msp{page 95}
of the boundary of the 2-simplex

\begin{center}
\begin{tikzpicture}[baseline=(b.base)]
\node (b) at (0.95,0.3) {};
\draw (0,0) -- (1.9,0) -- (0.95,0.95) -- cycle;
\fill (0,0) circle (2.2pt);
\fill (1.9,0) circle (2.2pt);
\fill (0.95,0.95) circle (2.2pt);
\end{tikzpicture}
\qquad\qquad\qquad or \qquad\qquad\qquad
\begin{tikzpicture}[baseline=(c.base)]
\node (c) at (0,-0.1) {};
\draw (0,0) circle (0.48);
\fill (90:0.48) circle (2.2pt);
\fill (210:0.48) circle (2.2pt);
\fill (330:0.48) circle (2.2pt);
\end{tikzpicture}
\qquad\quad ,
\end{center}

barycentric subdivision
\printedpage{48}
% head: A. Grothendieck : Dérivateurs, Ch. XVI, Édition des Universités de Montpellier II et Paris VII
% foot: Version 18 mars 2019

\begin{center}
\begin{tikzpicture}[baseline=(b.base)]
\node (b) at (0.95,0.3) {};
\draw (0,0) -- (0.95,0.95) -- (1.9,0);
\fill (0,0) circle (2.2pt);
\fill (0.95,0) circle (2.2pt);
\fill (1.9,0) circle (2.2pt);
\fill (0.475,0.475) circle (2.2pt);
\fill (1.425,0.475) circle (2.2pt);
\fill (0.95,0.95) circle (2.2pt);
\end{tikzpicture}
\qquad\qquad\qquad or \qquad\qquad\qquad
\begin{tikzpicture}[baseline=(c.base)]
\node (c) at (0,-0.1) {};
\draw (0,0) circle (0.48);
\fill (30:0.48) circle (2.2pt);
\fill (90:0.48) circle (2.2pt);
\fill (150:0.48) circle (2.2pt);
\fill (210:0.48) circle (2.2pt);
\fill (270:0.48) circle (2.2pt);
\fill (330:0.48) circle (2.2pt);
\end{tikzpicture}
\end{center}

(hexagon), that for every increasing map $g : I \longrightarrow \mathrm{sd}(I)$, its topological realization is homotopic to a constant map, so $g$ cannot be a homotopism in Cat (for this would imply that its topological realization is a homotopism).

This therefore shows that the notion of homotopism in Cat is \textit{strictly finer} than that for the topological realizations, which once again restores substance to the hope that Whitehead's invariant could be defined as a homotopy invariant on Cat, or at least on a suitable full subcategory, such as that of ordered sets (possibly restricted to be finite), or of simplicial mock-ups. And even, that the set of Wh-homotopy types (simple homotopy types) can be
\msp{page 96}
canonically embedded in a suitable $\mathrm{Hot}_W$, where $W$ would be a basic localizer strictly finer than $\mathbf{W}_\infty$, e.g. (who knows?) $\tilde{\mathbf{W}}_\infty$. The first piece of work one would like to do is to construct such a $W$ \textit{ad hoc}, so that Whitehead's invariant can be written as a functor on $\mathrm{Hot}_{\bullet W}$. To begin with, one would already have to make sure that this invariant can be viewed as an invariant of homotopy type, \textit{i.e.} that it is interpreted as a functor on $\mathrm{Hot}_{\bullet}\mathrm{Cat} \simeq \mathrm{Cat}(\,\underbrace{\text{homotop}}_{\substack{\text{set of}\\ \text{homotopisms}}}\,)^{-1}$. This work will perhaps also suggest the definition of ``higher Whitehead invariants'' (involving the $\pi_i$, not only the $\pi_1$), which could perhaps serve (with the one of dimension 1) to describe the hypothetical localizer $W$ (``of Whitehead'').

\msp{page 97}

\section*{7\quad Comments on the axioms}

7.1. The main reason for being of a basic localizer $W$ is to define a theory of homotopy’, and more precisely a prederivator $\mathrm{HOT}_W$, of which one hopes that it will even be a derivator, or that it will at least satisfy some of the axioms of derivators. The axioms Loc (1) to Loc (9) are for the most part oriented toward this objective, except that Loc (1, 2, 3) constitutes an absolute ``bare minimum'', to be able to begin working with $W$. The main questions which arise seem to me to be the following: % [sic] d'homotopie’ -> d'homotopie (or “théorie d'homotopie”), closing quote without opening quote
\begin{enumerate}[label={}, leftmargin=2.5em]
\item[1)] When is Cat equipped with $(W, \mathrm{Cof}_W)$, resp. with $(W, \mathrm{Fib}_W)$, a category with cofibrations resp. with fibrations, so that the results of K. S. \textsc{Brown} are applicable?
\item[2)] Under what conditions, on $W$ and on $f$, do $f_!$ and $f_*$ exist? And when do we have, in addition, exactness of cocartesian resp. cartesian squares?

\msp{page 98}
\item[3)] More generally, in terms of the axioms imposed on $W$, which are the axioms Der (\ldots) that are satisfied for the prederivator $\mathrm{HOT}_W$? (Which could also depend on the choice of the domain Diag $\subset$ Cat. The prederivator is defined a priori on the whole of Cat). In particular, % [sic] quelles sont les axiomes -> quels sont les axiomes, agreement; [sic] que sont vérifiés -> qui sont vérifiés, relative pronoun
\end{enumerate}

\end{document}
